\chapter{Swallowed substance}
\index{Swallowed substance}

\medskip
\noindent\textit{Educational reference; always seek professional advice for individual care.}

\section*{Definition and Overview}
"Swallowed substance" covers two related situations: swallowing something that was never meant to be swallowed, and swallowing too much of something that is normally harmless. The first group, foreign body ingestion, includes coins, buttons, small toys, bones, dentures, and the particularly dangerous button batteries and magnets. The second group, poisoning or overdose, includes medicines taken accidentally or deliberately, household cleaning products, alcohol, pesticides, plants, and toxic berries. Children, especially those between six months and three years, explore the world with their mouths and account for most accidental ingestions, while deliberate ingestion of medicines or poisons occurs at all ages and may be a sign of significant distress. Most swallowed objects pass through the digestive system without harm, and many poisonings are minor, but a small minority are genuinely dangerous, and the outcome often depends on how quickly help is sought. The most important steps for anyone facing this situation are not to panic and not to wait: advice from a poison information service or a doctor should be obtained promptly, and certain objects, especially button batteries, require emergency assessment. This chapter explains which swallowed substances are dangerous, what symptoms to watch for, how these situations are assessed and treated, and how accidental ingestion is prevented, in plain language for parents, carers, and anyone who witnesses an ingestion. Clinicians always prefer a phone call too early to an emergency too late.

\section*{Epidemiology}
Ingestions are extremely common. In many countries and comparable countries, calls to poisons information centres number in the hundreds of thousands each year, and children aged one to three years are involved in the majority. Coins are the commonest foreign body swallowed by children, and most pass naturally. The objects that cause the most serious harm are button batteries and multiple magnets, and although they are swallowed less often than coins, they are responsible for a disproportionate share of severe injuries and deaths. Accidental poisoning in children most often involves medicines kept in handbags, on benches, or in pill organisers, and household products such as dishwasher tablets, bleach, and detergents. Deliberate overdose is different: adolescents and adults, often with mental health difficulties, and increasingly involving paracetamol, which is freely available and a common cause of serious liver injury. The elderly are a further group at risk, because of confusion, poor eyesight, or mistakes with multiple medicines. Poisoning rates in many countries have fallen over recent decades thanks to child-resistant packaging and education, but cases still occur regularly, and the substances involved change over time as new products reach the market.

\section*{Aetiology and Pathophysiology}
The consequences of an ingestion depend on what was swallowed, how much, and how harmful it is. A swallowed object is handled in one of three ways: it may pass harmlessly through the digestive tract, become stuck along the way, or damage the tissues it contacts.

\begin{itemize}
    \item \textbf{Passage:} most small, smooth objects such as coins pass through the stomach and bowel and are passed in the stool within days, without harm
    \item \textbf{Impaction:} objects can lodge in the oesophagus, most commonly at the level of the voice box or lower oesophagus, causing pain, drooling, and difficulty swallowing
    \item \textbf{Button batteries:} a disc battery stuck in the oesophagus leaks an alkaline solution that burns through the wall of the oesophagus within hours, with the capacity to cause severe bleeding, perforation, and death
    \item \textbf{Magnets:} two or more magnets swallowed at different times can attract each other across the walls of the bowel, pinching and perforating the gut
    \item \textbf{Corrosive chemicals:} acids and alkalis such as bleach, oven cleaner, and dishwasher tablets burn the mouth, throat, and stomach, and can perforate the oesophagus
    \item \textbf{Poisonous medicines:} drugs such as paracetamol damage the liver, iron and some antidepressants damage various organs, and sedatives and alcohol suppress breathing
    \item \textbf{Alcohol:} even small amounts can cause dangerous low blood sugar and breathing suppression in small children
    \item \textbf{Plants and berries:} a few common local plants are poisonous, although most cause only nausea and vomiting
\end{itemize}

The speed of the process varies greatly: a caustic burn and a button battery damage tissue within hours, paracetamol takes a day or more to show liver injury, and most coins cause problems only if they become stuck. This is why the identity of the substance and the timing of ingestion are the two facts clinicians most need to know.

\section*{Clinical Features}
Symptoms depend on what was swallowed, but many children with an ingestion have no symptoms at all, and the event may be discovered only because a parent noticed the substance was missing.

\begin{itemize}
    \item \textbf{No symptoms:} the commonest situation, especially with coins and small objects that pass straight through
    \item \textbf{Drooling and refusing to swallow:} suggests something is stuck in the oesophagus
    \item \textbf{Pain in the throat or chest:} can indicate an impacted object or a burn from a corrosive substance
    \item \textbf{Vomiting, sometimes with blood:} occurs with oesophageal irritation, corrosive injury, or some poisons
    \item \textbf{Coughing, choking, or gagging:} may mean the object has entered the airway rather than the food pipe
    \item \textbf{Abdominal pain:} with large objects, magnet chains, or corrosive injury
    \item \textbf{Nausea, drowsiness, or odd behaviour:} common with medicines and alcohol, and the more alarming the behaviour, the more urgent the assessment
    \item \textbf{Burning pain in the mouth:} typical of corrosive household products, often with visible burns to the lips and tongue
    \item \textbf{Blood in vomit or black stools:} signs of bleeding from the stomach or bowel
\end{itemize}

\section*{Differential Diagnosis}
The challenge with a suspected ingestion is that symptoms are non-specific, and clinicians must work out whether the child is unwell from the swallowed item or from another illness.

\begin{itemize}
    \item \textbf{Gastroenteritis:} vomiting and abdominal pain are common in childhood infections and may be mistaken for the effects of an ingestion
    \item \textbf{Croup or a respiratory infection:} cough and noisy breathing can be confused with an object entering the airway, although true choking is sudden
    \item \textbf{Severe tonsillitis:} drooling and refusing to swallow can be caused by a sore throat rather than an impacted object
    \item \textbf{Constipation:} abdominal discomfort may be blamed on a swallowed object that was never there, or vice versa
    \item \textbf{Appendicitis:} central abdominal pain in a child may be the first sign of appendicitis rather than the result of an ingestion
    \item \textbf{Mental health crises:} in older children and adults, deliberate ingestion is often a symptom of underlying distress rather than of poisoning itself
\end{itemize}

\section*{When to Seek Medical Care}
A GP or the Poisons Information Centre (in many countries, 13 11 26) should be contacted for any significant ingestion, even if the person appears well. Do not wait for symptoms before seeking advice about a swallowed button battery, magnet, or corrosive substance. Do not induce vomiting, and do not give anything to eat or drink unless a clinician advises it.

\textbf{Contact your local emergency services for:}
\begin{itemize}
    \item Any swallowed button battery or multiple magnets, or any child with drooling, trouble breathing, or severe pain
    \item Difficulty breathing, choking, or a cough that continues after the initial event
    \item Collapse, seizures, or altered level of consciousness
    \item Vomiting blood, or severe abdominal pain with a rigid or swollen belly
    \item A known deliberate overdose, even if the person appears well, because some effects are delayed
    \item Any person who is confused, suicidal, or unable to be roused
\end{itemize}

\section*{Diagnosis and Investigations}
Assessment begins with the history: what was swallowed, how much, when, and what symptoms have developed. Clinicians may ask to see the container or the remaining object to confirm the identity and dose.

\begin{itemize}
    \item \textbf{Poisons Information Centre advice:} the first step for most ingestions; specialists calculate the risk based on the substance, amount, and weight of the person, and advise whether observation, a GP visit, or emergency care is needed
    \item \textbf{X-ray:} coins, button batteries, and many other objects are visible on a plain X-ray of the chest and abdomen
    \item \textbf{Endoscopy:} a camera examination of the oesophagus and stomach, used to remove an impacted object or battery and to assess corrosive burns
    \item \textbf{Blood tests:} paracetamol levels, liver function, and electrolyte tests are used for overdoses and are timed according to the drug ingested
    \item \textbf{Observation:} many children are watched in hospital for a period to ensure the object passes
    \item \textbf{Electrocardiogram:} some medications affect the heart rhythm, so an ECG may be needed after certain overdoses
    \item \textbf{Psychiatric assessment:} after a deliberate overdose, a mental health assessment is an essential part of care, in hospital and then in follow-up
\end{itemize}

\section*{Management}
Treatment ranges from simple reassurance to emergency surgery, and the plan depends entirely on what was swallowed.

\begin{itemize}
    \item \textbf{Watch and wait:} for small smooth objects like coins and most buttons, the standard advice is to let the object pass naturally, checking the stools until it appears, usually within three to five days
    \item \textbf{Urgent removal:} button batteries in the oesophagus are a surgical emergency and must be removed as fast as possible by endoscopy, because of the risk of catastrophic burns
    \item \textbf{Endoscopic removal:} objects stuck in the oesophagus, or sharp objects and magnets, are removed under anaesthesia
    \item \textbf{Surgery:} rarely needed, but required if an object or a chain of magnets has perforated the bowel
    \item \textbf{Antidotes:} specific medicines reverse specific poisons, such as acetylcysteine for paracetamol and naloxone for opioid overdoses
    \item \textbf{Supportive care:} intravenous fluids, medicines for symptoms, and observation are the mainstay for most poisonings
    \item \textbf{Treatment of burns:} corrosive injuries need careful assessment, often with endoscopy, and sometimes antibiotics and nutritional support while the oesophagus heals
    \item \textbf{Mental health care:} for deliberate overdose, safety planning, counselling, and referral to community mental health services are essential and non-judgemental
\end{itemize}

\section*{Complications}
\begin{itemize}
    \item \textbf{Oesophageal injury:} a stuck object or corrosive burn can cause ulceration, stricture (narrowing), or perforation
    \item \textbf{Perforation of the bowel:} the most serious complication of magnet and sharp-object ingestion, requiring surgery
    \item \textbf{Liver damage:} paracetamol overdose is a leading cause of acute liver failure and may require intensive care or transplantation
    \item \textbf{Bleeding:} ulceration from batteries or corrosives can bleed, occasionally severely
    \item \textbf{Aspiration:} vomited material or a partly swallowed object can enter the lungs, causing pneumonia
    \item \textbf{Long-term psychological harm:} in deliberate overdose, the mental health crisis that led to the ingestion is itself the most important complication to address
\end{itemize}

\section*{Prognosis and Outcomes}
The overwhelming majority of people who swallow a foreign body or a poison recover completely. Most foreign bodies pass without treatment, most poisonings are managed with observation and supportive care, and with prompt removal the outlook after a button battery or impacted object is good. The exceptions are a minority of cases: a button battery left in the oesophagus can be fatal or cause permanent damage, delayed treatment of a corrosive burn can leave a narrowed oesophagus, and untreated paracetamol overdose can destroy the liver. These outcomes are all time-sensitive, which is why early, appropriate help is the strongest predictor of a good result. For people who have taken an overdose deliberately, the medical outcome is usually good, but the underlying distress remains, and the outlook for their wellbeing depends on compassionate mental health support. Families are encouraged to seek help for a child's accidental ingestion without shame or blame: accidents happen to every family, and clinicians respond quickly and non-judgementally.

\section*{Prevention and Self-Care}
\begin{itemize}
    \item Store medicines, cleaning products, batteries, and small objects out of sight and out of reach, in locked or child-resistant cabinets
    \item Keep button batteries and loose magnets away from children, and check that toys, remote controls, and musical cards have secure battery compartments
    \item Keep the Poisons Information Centre number (13 11 26) and 000 where you can see them in an emergency
    \item Keep medicines in their original, labelled containers and never call tablets "lollies"
    \item Dispose of old medicines through a pharmacy return scheme rather than in household bins
    \item Never induce vomiting after an ingestion unless a professional tells you to
    \item If you or someone you know has swallowed something deliberately, seek help immediately: emergency services, your GP, a a local crisis service, a a local mental-health support service, and for young people a a local youth crisis service, and for Aboriginal and Torres Strait Islander people a a local culturally appropriate crisis service
\end{itemize}

\section*{Sources and Further Reading}
The following organisations provide reliable information about poisons, foreign body ingestion, and overdose prevention.

\begin{itemize}
    \item healthdirect Australia. \textit{Poisoning and first aid for swallowed substances.} https://www.healthdirect.gov.au/
    \item Poisons Information Centre Australia. \textit{What to do in a poisoning emergency.} https://www.poisonsinfo.nsw.gov.au/
    \item NHS. \textit{Swallowed objects and accidental poisoning in children.} https://www.nhs.uk/
    \item Mayo Clinic. \textit{First aid: poisoning and swallowed foreign objects.} https://www.mayoclinic.org/
    \item Australian Department of Health and Aged Care. \textit{Medicine safety and poison prevention.} https://www.health.gov.au/
    \item World Health Organization (WHO). \textit{Poison prevention and chemical safety.} https://www.who.int/
\end{itemize}

\clearpage
